\section{Further questions}\label{sec:questions}

The mass counterexample and the cancellation theorems concern the same
positive phase equation. They show that acyclic cancellation need not
come from trace-norm mass dissipation. For an invertible differential,
the contracting quotient $A_jD^{-1}$ supplies the missing control from
every initial metric. With a surviving source, estimates using only the
nonsingular cohomology block close the argument at large initial acyclic
scale. They allow arbitrary surviving tuples, unequal dimensions and
fixed upper-half-plane charges, uniformly in arrow count, stabilizers
and acyclic condition numbers.

Cancellation and relaxation have different hypotheses. For equal surviving
dimensions and charges $-1+\ii$, $1+\ii$, a finite infimum of the
cohomology Kempf--Ness function gives relaxation from nonbalanced initial
metrics. Its quantitative energy bound retains the initial excess
$\mathcal K(0)-\inf\mathcal K$. An arrow-norm bound alone need not bound
this excess: for a scalar one-arrow representation $B=b\ne0$ at standard
metrics, it is $|b|^2-1-\log|b|^2$, which diverges as $b\to0$.
Lower boundedness, and not attainment, is the analytic hypothesis used
in the proof. For arbitrary surviving representations that hypothesis
can fail even though large-scale cancellation still holds.

Removing the large-initial-scale assumption remains open for the mixed
theorem. The inequality
$s'\leq(-a+4\pi q/(b_0\lambda))s$ requires control of $q/\lambda$
before it becomes contractive; monotonicity of $\lambda$ alone supplies
no such bound. It is also natural to ask whether the mechanism extends
to longer complexes while retaining independence from positive
surviving orbit eigenvalues.

The integrated moment-map estimate gives no uniform final relaxation
rate, normalized surviving-metric limit, or parameter derivatives of
actual trajectories. These questions require additional control of the
surviving dynamics. The static results allow arbitrary fixed invertible
acyclic differentials but do not differentiate that differential.
For the explicit scalar-scale family, mass converges in $C^2$ and every
positive H\"older seminorm of its Hessian diverges. Optimal general
second-derivative estimates on the full arrow-coordinate space remain open.

The explicit primitive separates metric-direction exactness from
transverse positivity. Its negative ambient Levi direction excludes
one direct positive-potential construction. The limiting cohomology
metric and restricted finite-scale example leave open a general
quotient-only construction with its own normalization and compatibility
laws. Extensions to bundle complexes and singular geometric limits
require additional elliptic analysis; none follows from these
finite-dimensional results.
